\section{Conditional-count estimates}\label{v9:count:appendix}
\label{ms:fourier:local}
This appendix supplies the analytic estimates needed after confinement.
Its statements are conditional on their displayed support, probability,
count and sign hypotheses. In particular, the formulas requiring $p<1/2$
do not cover the entire acquired range $[.25,.53]$. Neither reflection nor
the half-amplitude result alone fills that gap. No complete numerical
comparison cover is asserted here.

\subsection{The comparison and its domain}\label{ms:fourier:cut49:shortloc:section}

A law concentrated near two points has a natural integer count and a small
centered error. We compare it with its own Bernoulli count. The same exact
identity handles every noise regime. There are two useful estimates of its
right side: discard all but the nearest outgoing count, or retain all counts
through a random cutoff. The estimates share their normalization, moments,
and tail bounds; only the final scalar witness changes.

\paragraph{The identity and sufficient inequalities}

Let $0<p\le1/2$, $q=1-p$, $v=pq$, and let
\[
 Y=B+U,\qquad \Pp(B=1)=p,\qquad \E(U\mid B)=0.
\]
The two actual conditional laws $U_0,U_1$ obey
\begin{equation}\label{ms:fourier:cut49:shortloc:support}
 |U_j|\le A,\quad U_0\le a_0,\quad U_1\ge-a_1,\quad
 a_0,a_1\ge0,\quad d=a_0+a_1\le d_0<\min(p,q),\quad d_0\le A.
\end{equation}
Write $s_j=\E U_j^2$, $s=qs_0+ps_1$, $\lambda=ns$,
$\ell=(v+s)^{-1/2}$, $\theta=s/(v+s)$ and
$\rho=\E|\ell(Y-p)|^3$. At the actual threshold choose
\[
 x=np+H/\ell=k+u,\qquad k\in\{1,\ldots,n-1\},\quad |u|\le\tfrac12,
 \qquad \delta=k-np=H/\ell-u.
\]
Either nearest integer is allowed at a tie. Put
\[
 \pi_j=\binom njp^jq^{n-j},\quad F_B(j)=\Pp(\operatorname{Bin}(n,p)\le j),
 \qquad \rho_B=(1-2v)/\sqrt v,
\]
with $\pi_j=0$ outside $[0,n]$. Let $W_j$ be an independent sum of $j$
copies of $U_1$ and $n-j$ copies of $U_0$. The signed scores are
\[
 C=\frac{\sqrt n}{\rho}\{\Pp(\textstyle\sum Y_i\le x)-\Phi(H/\sqrt n)\},
 \qquad C_B=\frac{\sqrt n}{\rho_B}\{F_B(k)-\Phi(\delta/\sqrt{nv})\}.
\]
The fixed target is $c=\CE$, with $c\ge c_*>0$; one may take
$c_*=0.4$. The actual ratio $C$ is used only to name the excess
$a=\rho(C-c)$. It is not the coefficient in the moment correction.
Define
\begin{align}
 f(u)&=\frac{\Phi((\delta+u)/\sqrt{n(v+s)})-
                    \Phi(\delta/\sqrt{n(v+s)})}{\pi_k},\label{ms:fourier:cut49:shortloc:f}\\
 Q&=\Pp(\textstyle\sum Y_i\le k+u)-F_B(k),\nonumber\\
 \mathcal R&=\sqrt n\{\Phi(\delta/\sqrt{nv})-
                 \Phi(\delta/\sqrt{n(v+s)})\}+c(\rho_B-\rho).\nonumber
\end{align}
Equation~\eqref{ms:master-identity}, with $V_j=W_j$ and the same
interior count $k$, gives the exact identity
\begin{equation}\label{ms:fourier:cut49:shortloc:master}
 a+\rho_B(c-C_B)=\sqrt n\pi_k\{Q/\pi_k-f(u)\}+\mathcal R.
\end{equation}
In particular, the phase $u$, the sign of $H$, and both conditional laws
refer to the same marginal throughout.

\begin{theorem}[Parameterized local comparison]\label{ms:fourier:cut49:shortloc:theorem}
Fix $n\ge N$, $\lambda>0$ and a closed parameter cell on which
\eqref{ms:fourier:cut49:shortloc:support} holds. Suppose $\sqrt n\pi_k\ge m>0$.
Either of the following witnesses proves $a+\rho_B(c-C_B)<-\varepsilon\lambda$.
\begin{enumerate}
\item \emph{Nearest-count witness.} Valid bounds $O,T,R_0,\kappa$ satisfy
\[
 \Pp(W_k>u)+f(u)\ge O,\quad
 \sum_{i\ge1}\Pp(W_{k+i}\le u-i)\le T,\quad
 \pi_{k+i}/\pi_k\le R_0,\quad \mathcal R\le\kappa\lambda/n.
\]
Only legal counts occur in the sum; adding positive upper bounds for
other counts is allowed. If $D=O-R_0T>0$ and
\begin{equation}\label{ms:fourier:cut49:shortloc:directtest}
             mD-\kappa\lambda/n>\varepsilon\lambda,
\end{equation}
the conclusion follows.
\item \emph{Whole-count witness.} Bounds $B(\lambda)<0$ and $I(\lambda)\ge0$
satisfy
\[
 a+\rho_B(c-C_B)\le\sqrt n\pi_k\lambda B(\lambda)
                                      +\lambda I(\lambda)/n.
\]
It is enough that
\begin{equation}\label{ms:fourier:cut49:shortloc:whole-test}
                  mB(\lambda)+I(\lambda)/N<-\varepsilon.
\end{equation}
\end{enumerate}
The witnesses may differ from cell to cell. A finite closed cover and a
separate treatment of $\lambda=0$ therefore suffice for a complete local
comparison.
\end{theorem}
\begin{proof}
The signed count identity \eqref{ms:count-difference}, with $V_j=W_j$, is
\begin{equation}\label{ms:fourier:cut49:shortloc:count}
 Q=-\sum_{i=0}^{k}\pi_{k-i}\Pp(W_{k-i}>u+i)
       +\sum_{i=1}^{n-k}\pi_{k+i}\Pp(W_{k+i}\le u-i).
\end{equation}
Apply Proposition~\ref{ms:local-witnesses} with $m_0=m$.
The legal counts, closed incoming events and strict outgoing events are
identical. In the second alternative the present hypothesis $I\ge0$
makes its $\max(0,I)/N$ exactly $I/N$, yielding
\eqref{ms:fourier:cut49:shortloc:whole-test}.
\end{proof}

\paragraph{Binomial normalization and the variance correction}

Here is an explicit finite recipe for the constants required above.
Let $p\in[p_-,p_+]\subset(0,1/2)$, $n\ge N$, $\lambda\le\Lambda$,
$\theta\le\Theta$, and $|\delta|\le\Delta$, where $\Delta\ge1$ also includes
every binomial mode. Choose $0<r_-<p_- -\Delta/N$ and
$p_++\Delta/N<r_+<1/2$. Let $v_*>0$ be a lower bound for $p(1-p)$.
Choose positive $E,L,R$ bounding, on $[r_-,r_+]$,
\[
 |D'''(r\Vert p)|=\frac{1-2r}{r^2(1-r)^2},\qquad
 |(\log(r(1-r)))'|=\frac{1-2r}{r(1-r)},\qquad
 \frac1{12r(1-r)},
\]
respectively, where $D(r\Vert p)=r\log(r/p)+(1-r)\log((1-r)/q)$.
Then, with $Z(t)=Lt/2+R+Et^3/(6N)$,
\begin{equation}\label{ms:fourier:cut49:shortloc:mass}
 \sqrt n\pi_k=\frac{\phi(0)}{\sqrt v}
       e^{-\delta^2/(2nv)+\zeta},\qquad n|\zeta|<Z(|\delta|).
\end{equation}
For completeness, the finite Stirling formula used here is
\[
 \log(j!)=(j+\tfrac12)\log j-j+\tfrac12\log(2\pi)+r_j,
 \qquad 0<r_j<1/(12j).
\]
If $a_j=\log(j!)-(j+1/2)\log j+j$ and $z=(2j+1)^{-1}$, the integrated
geometric series gives
$0<a_j-a_{j+1}=\sum_{i\ge1}z^{2i}/(2i+1)<1/[12j(j+1)]$.
Summing gives the remainder relative to $\lim a_j$. The integration-by-parts
recursion for $\int_0^{\pi/2}\sin^j t\,dt$ and the inequalities between three
successive integrals give Wallis's product; its finite product equals
$e^{4a_j-2a_{2j}}j/[2(2j+1)]$ and tends to $\pi/2$. Thus
$\lim a_j=\tfrac12\log(2\pi)$. Applying this formula to $n!,k!,(n-k)!$
and Taylor-expanding $D$ and the logarithmic prefactor proves
\eqref{ms:fourier:cut49:shortloc:mass}; the factorial error is at most $1/[12nr(1-r)]$.

Consequently a valid mass lower bound is any positive
\begin{equation}\label{ms:fourier:cut49:shortloc:m}
 m<\frac{0.39886}{\sqrt{p_+(1-p_+)}}
          \left[1-\frac{\Delta^2/(2v_*)+Z(\Delta)}N\right].
\end{equation}
The same formula at a mode proves $\sqrt n\max_j\pi_j<1$ whenever
$v_*[1-Z(\Delta)/N]^2>0.4^2$ and $Z(\Delta)<N$.
The rational bounds $0.39886<\phi(0)<0.4$ follow, for example, by the
alternating arctangent series in Machin's identity.

Differentiating \eqref{ms:fourier:cut49:shortloc:f} gives
\begin{equation}\label{ms:fourier:cut49:shortloc:logder}
 \log f'=-\tfrac12\log(1+s/v)+\delta^2/(2nv)
                   -(\delta+u)^2/(2n(v+s))-\zeta.
\end{equation}
Thus $n|\log f'|\le B_0\lambda+P(D,a)$ on $|\delta|\le D\le\Delta$,
$|u|\le a\le1/2$, where
\[
 B_0>\frac1{2v_*}+\frac{\Delta^2}{2Nv_*^2},\qquad
 P(D,a)=\frac{2Da+a^2}{2v_*}+Z(D).
\]
Choose $0<b<1$ above $[B_0\Lambda+P(\Delta,1/2)]/N$ and put
$F=(1-b)^{-1}$. The elementary estimates $e^z\le(1-z)^{-1}$ and
$1-e^{-z}\le z$ imply
\begin{equation}\label{ms:fourier:cut49:shortloc:derivative}
 |f'-1|\le(B_1\lambda+P_1)/n,\qquad B_1=FB_0,\quad P_1=FP(D,a).
\end{equation}
This supplies both a derivative interval $l_f\le f'\le h_f$ and the
finer phase-error bound $|u-f(u)|\le a(B_1\lambda+P_1)/n$.

The exact count variance is
\begin{equation}\label{ms:fourier:cut49:shortloc:variance}
 t_j=\E W_j^2=\lambda+(j-np)(s_1-s_0),\qquad
 |s_1-s_0|\le s/p\le c_v s,\quad c_v\ge1/p_-.
\end{equation}
It gives $\omega_-\lambda\le t_k\le\omega_+\lambda$ with
$\omega_\pm=1\pm c_vD/N$. All uses below require the checked condition
$\omega_->0$. Also
$t_{k+i}\le iV_*\lambda$ for $V_*=1+c_v(\Delta+1)/N$.
The binomial recursion gives, for every legal $i\ge1$,
\begin{equation}\label{ms:fourier:cut49:shortloc:ratio}
 \frac{\pi_{k+i}}{\pi_k}\le R_0,
 \qquad R_0>1+\frac{p_++\Delta-1}
 {Nv_*[1-(\Delta-1)/(Np_-)]},
\end{equation}
provided the denominator is positive and $p_++\Delta<2$.
Indeed the first ratio has this bound and all subsequent rising factors
are at most one. No count outside $[0,n]$ is used in this recursion.

Finally \eqref{ms:fourier:cut49:shortloc:support} fixes the signs inside the absolute cubic.
Expansion and centering give
\[
 \E|Y-p|^3=v(1-2v)+3pq(s_0+s_1)
                   +p\E U_1^3-q\E U_0^3
 \ge v(1-2v)+(3p-d)s.
\]
The inequality $(1+z)^{3/2}\le1+3z/2+3z^2/8$ for $z\ge0$ yields
\begin{equation}\label{ms:fourier:cut49:shortloc:cubic}
 (v+s)^{3/2}(\rho-\rho_B)\ge s
     \left(6p-3p^2-\tfrac32-d-\frac{3(1-2v)s}{8v}\right).
\end{equation}
Let $Q_*>0$ be a lower bound for the bracket and
$m_c\le c_*/[p_+(1-p_+)+\Lambda/N]^{3/2}$.
The normal variance correction is nonpositive for $\delta\le0$;
for $\delta>0$ its derivative in $s$ is at most $0.4\delta/(2v^{3/2})$.
Therefore
\begin{equation}\label{ms:fourier:cut49:shortloc:repair}
 \mathcal R\le(P_*D_+-m_cQ_*)\lambda/n,
 \qquad P_*\ge\frac{0.4}{2v_*^{3/2}},\quad D_+\ge\delta_+.
\end{equation}
One may take $D_+=D$, but retaining the positive part is often decisive.
All bounds concern $v+s$, which can exceed $1/4$.

\paragraph{Two elementary tools for probabilities}

\begin{lemma}[A shifted quartic]\label{ms:fourier:cut49:shortloc:quartic}
Suppose $W$ is a sum of independent centered variables bounded in absolute
value by $A$, and $t=\E W^2$. Let $K_0\ge(9+6\sqrt3)/4$ and $a>0$.
For $r\ge0$ and $\sigma\in\{-1,1\}$,
\begin{equation}\label{ms:fourier:cut49:shortloc:quarticformula}
 \Pp(W+\sigma r>0)\ge
 \frac{(3a^2-A^2)t-3t^2-4Art-6r^2t-r^4
                      +3a^2r^2+2\sigma a^3r}{K_0a^4}.
\end{equation}
\end{lemma}
\begin{proof}
The polynomial $h_a(z)=(-z^4+3a^2z^2+2a^3z)/(K_0a^4)$ is below
$\mathbf1_{\{z>0\}}$: it is nonpositive for $z\le0$, and its maximum for
$z>0$ is at $z/a=(1+\sqrt3)/2$ and is at most one. Independence gives
$|\E W^3|\le At$ and $\E W^4\le A^2t+3t^2$.
Expand $\E h_a(W+\sigma r)$ and apply these two inequalities.
The value at the threshold is zero, so atoms require no limiting argument.
\end{proof}

\begin{lemma}[One common tail bound]\label{ms:fourier:cut49:shortloc:tails}
If independent centered $X_i\le A$ have total variance at most $V$, then,
for $b>0$, their sum has upper tail bounded by the minimum of one and
\begin{equation}\label{ms:fourier:cut49:shortloc:tailmin}
 \begin{gathered}
 \frac{V}{V+b^2},\qquad
 \mathcal B(A,V,b)=e^{b/A}\left(\frac{V}{V+Ab}\right)^{b/A+V/A^2},\qquad
 \frac{\E(P_\mu-j)_+^3}{(b/A+\mu-j)^3},\\
 \mu=V/A^2,\qquad 0\le j<b/A+\mu.
 \end{gathered}
\end{equation}
where $P_\mu$ is Poisson. The same holds for a lower tail when $|X_i|\le A$.
All bounds apply to closed events, hence also to strict events.
\end{lemma}
\begin{proof}
Cantelli's bound follows by applying Markov to $(\sum X_i+V/b)^2$.
Taylor's integral remainder gives
$\E e^{zX_i}\le\exp\{\E X_i^2(e^{zA}-1-zA)/A^2\}$
for $z\ge0$. Multiply and use $z=A^{-1}\log(1+Ab/V)$ to get $\mathcal B$.
The zero-variance case is immediate.

For the third bound, smooth $f(x)=(x-t)_+^3$ by a nonnegative smooth
compactly supported density. Then $f'',f''',f''''\ge0$. For
$\mathcal Af(x)=f(x+A)-f(x)-Af'(x)$, the Taylor quotient
\[
 \frac{f(x+y)-f(x)-yf'(x)}{y^2}
       =\int_0^1(1-r)f''(x+ry)\,dr
\]
is nondecreasing in $y$. A centered $X\le A$ of variance $v$ consequently
satisfies $\E f(x+X)\le f(x)+(v/A^2)\mathcal Af(x)$.
For $\mathcal P_zf(x)=\E f(x+A(P_z-z))$,
$\partial_z\mathcal P_zf=\mathcal P_z\mathcal Af$ and
\[
 \partial_z^2\mathcal P_zf=\mathcal P_z\mathcal A^2f\ge0,\qquad
 \mathcal A^2f(x)=\int_0^A\!\int_0^A(A-r)(A-w)f''''(x+r+w)\,dr\,dw.
\]
Thus $\mathcal P_{v/A^2}f\ge f+(v/A^2)\mathcal Af$. Replace the independent
variables successively; translations and expectations preserve the derivative
signs and Poisson rates add. Increasing the total variance to $V$ is permitted
by Jensen after adding independent centered Poisson noise. Removing the
smoothing is justified by bounded actual upper supports and finite Poisson
third moments. Markov at $t=A(j-\mu)$ proves the displayed bound.
\end{proof}

These bounds have finite rational evaluations. For example,
\begin{align}\label{ms:fourier:cut49:shortloc:poissonfinite}
 \E(P_\mu-j)_+^3={}&\mu^3+3(1-j)\mu^2+(1-3j+3j^2)\mu-j^3\\
 &+e^{-\mu}\sum_{h=0}^{\lfloor j\rfloor}\frac{\mu^h(j-h)^3}{h!}.\nonumber
\end{align}
Replacing $e^{-\mu}$ by $(\sum_{h=0}^{24}\mu^h/h!)^{-1}$ enlarges the
nonnegative correction. Alternatively sum the positive Poisson series
through 32; the remaining tail is at most its first term divided by
$1-\mu(34-j)^3/[34(33-j)^3]$, provided $j<33$ and that ratio is below one.
For Bennett, put $x=Ab/V$ and
$L=[(b/A+\mu)\log(1+x)-b/A]$. The positive series
$\log(1+x)=2\sum_{h\ge0}z^{2h+1}/(2h+1)$, $z=x/(2+x)$, gives lower
bounds for $L$; any positive lower exponential series gives an upper bound
for $e^{-L}$. Rounding a positive lower $L$ down before exponentiation is
also valid. These are choices of rational witnesses, not optimization claims.

\paragraph{The nearest-count witness, including vanishing noise}

For $u=-r\le0$, add $-h_fr$ to \eqref{ms:fourier:cut49:shortloc:quarticformula} with
$\sigma=1$; for $u=r\ge0$, add $l_fr$ with $\sigma=-1$.
Both bound $\Pp(W_k>u)+f(u)$ from below. For a cell
$t\in[t_-,t_+]$, $r\in[r_-,r_+]$, this expression is concave in $t$
(second derivative $-6/(K_0a^4)$), so only $t_-,t_+$ are needed.
At each endpoint, bound a monomial with positive coefficient at $r_-$
and one with negative coefficient at $r_+$. Any rational $a>0$ is a valid
witness. Other valid choices are
\begin{equation}\label{ms:fourier:cut49:shortloc:ownother}
 \frac{r_-^2}{t_++r_-^2}-h_fr_+\quad(u\le0),\qquad l_fr_-\quad(u\ge0).
\end{equation}
They follow from Cantelli and nonnegativity, respectively. Their maximum
or one selected witness is valid without an optimality proof.

For incoming counts use $t_{k+i}\le iV_*\lambda$. On $u\in[-r_+,-r_-]$
the first distance is $1+r_-$ and every later distance is at least $i$.
On $u\in[r_-,r_+]$ they are $1-r_+$ and $i(1-r_+/2)$ for $i\ge2$.
Let $b,b_o$ be these first and later distances. On $\lambda\in[L,U]$,
let $T_1$ be any bound from Lemma~\ref{ms:fourier:cut49:shortloc:tails} at $(A,V_*U,b)$.
Since $\mathcal B(A,iV,ib)=\mathcal B(A,V,b)^i$, all later counts total at most
$B^2/(1-B)$ for $B=\mathcal B(A,V_*U,b_o)<1$. Thus
\begin{equation}\label{ms:fourier:cut49:shortloc:midtest}
 D=O-R_0\{T_1+B^2/(1-B)\}>0,\qquad
 mD-\kappa U/N>\varepsilon U
\end{equation}
proves the entire cell when $\kappa\ge0$. This sign condition is required on each cell. The argument applies to
both central and left probability panels.

A cell touching $\lambda=0$ requires estimates proportional to $\lambda$;
a positive lower endpoint may not be substituted. The following rescaling
provides them. Put $r=zt$, $t_*=\omega_+U$ and $d_a=K_0a^4$.
The shifted quartic plus $br$, divided by $t$, becomes
\begin{equation}\label{ms:fourier:cut49:shortloc:scaledown}
 F_{\sigma,b}(z,t)=\frac{3a^2-A^2-3t}{d_a}
 +\left(\frac{2\sigma}{K_0a}+b-\frac{4At}{d_a}\right)z
 +\frac{3a^2t-6t^2}{d_a}z^2-\frac{t^3z^4}{d_a}.
\end{equation}
It extends to $t=0$ and is concave in $t$. On $0\le z\le z_0$, replacing
$z^4$ by $z_0^2z^2$ gives a lower quadratic at $t=t_*$; at $t=0$ it is
linear. Endpoints and a clipped quadratic vertex therefore give an exact
lower bound $J(\sigma,b,z_0)$ for the whole rectangle.

For negative phase choose $z_-,c_->0$ and verify
\begin{gather}\label{ms:fourier:cut49:shortloc:tinynegative}
 J(1,-h_f,z_-)>c_-,\quad
 \frac{z_-^2}{1+z_-^2t_*}-h_fz_->c_-,\quad
 \frac{2z_-}{(1+z_-^2t_*)^2}>h_f,\\
 \frac1{1+4t_*}-\frac{h_f}{2}>c_-t_*,\qquad z_-t_*<1/2.\nonumber
\end{gather}
The quartic handles $r\le z_-t$. Beyond that point use
$r^2/(t+r^2)-h_fr$. Its derivative at $r=z_-t$ is positive by the third
gate; its second derivative has just one sign change, so any later critical
point is a maximum. Its minimum is at $z_-t$ or $1/2$, covered by the other
gates. Hence the own count exceeds $c_-t$.
For positive phase choose $z_+,c_+>0$ with
\begin{equation}\label{ms:fourier:cut49:shortloc:tinypositive}
 J(-1,l_f,z_+)>c_+,\qquad l_fz_+>c_+,\qquad z_+t_*<1/2.
\end{equation}
The quartic handles $r\le z_+t$ and the normal term handles the rest.
The own count exceeds $c_+t$ and also exceeds $l_fr$.

For incoming tails proportional to $\lambda$, fix a Poisson anchor $j$.
A positive series term of index $h$ divided by
$\lambda(b/A+\mu-j)^3$ has logarithmic derivative, multiplied by $\lambda$,
\[
 h-1-\mu-\frac{3\mu}{b/A+\mu-j},\qquad \mu=V_*\lambda/A^2.
\]
If $b/A-j>0$ and this is nonnegative at $\lambda=U$ for the first positive
index, it is nonnegative for every term and all $0<\lambda\le U$.
Therefore the Poisson bound divided by $\lambda$ is largest at $U$.
Likewise $B^2/[\lambda(1-B)]$ is increasing whenever
$\lambda B'/B\ge1/2$. Differentiating the Bennett formula gives
\[
 \lambda B'/B=b/A-\mu\log(1+b/(A\mu));
\]
The function $g(\mu)=b/A-\mu\log(1+b/(A\mu))$ decreases because
$\log(1+x)>x/(1+x)$ for $x>0$. Moreover the logarithmic derivative of
$B^2/[\lambda(1-B)]$, multiplied by $\lambda$, is
$2g-1+gB/(1-B)\ge0$ when $g\ge1/2$. Thus checking the gate at $U$
proves the whole interval. A rational upper bound for the logarithm
verifies that gate. On cells using the fixed positive bound $l_fr_-$ instead,
unscaled endpoint tails suffice: after a positive bracket is checked, divide
it by $U$ to obtain a lower bound per $\lambda$. These arguments, followed by
\eqref{ms:fourier:cut49:shortloc:directtest}, prove the nearest-count comparison uniformly down to arbitrary positive
noise, including a deterministic conditional component.

\paragraph{The whole-count witness}

Assume $\zeta=1-d_0>1/2$. Couple all $2n$ conditional variables independently and
set $W_j=\sum_{i\le j}U_{1,i}+\sum_{i>j}U_{0,i}$, extending it constantly
outside $[0,n]$. Every increment of $j+W_j$ is at least $\zeta$. Define
\[
 K=\max\{j\in\mathbb Z:j+W_j\le k+u\},\qquad R=K-k.
\]
The coupling of Section~\ref{ms:counts}, with $V_j=W_j$, gives
\begin{equation}\label{ms:fourier:cut49:shortloc:cutoff}
                  \Pp(\textstyle\sum Y_i\le k+u)=\E F_B(k+R).
\end{equation}
Let $\chi=(n+1)p-k<2$ and $m_B=\max_j\pi_j$. For every integer $r$,
\begin{equation}\label{ms:fourier:cut49:shortloc:curvature}
 F_B(k+r)-F_B(k)-r\pi_k\le\frac{m_B}{(n+1)v}
 \left[(\chi-1)_+r_++\chi_+\binom{r_-}{2}+\binom{r_-}{3}\right].
\end{equation}
This is Lemma~\ref{ms:crossing-curvature}, with the same zero extension,
$\chi<2$ and interior count. Its proof requires only an increment lower
bound $\zeta>1/2$. Here $d\le d_0$ gives increments at least $1-d_0=\zeta$,
and $\zeta\le1$ follows from $d_0\ge0$.

Write $t=\E W_k^2$. Crossing monotonicity and Cauchy--Schwarz give
\begin{equation}\label{ms:fourier:cut49:shortloc:crossingmoments}
 \E R_+\le\frac{t}{(\zeta-u)^2},\quad
 \E\binom{R_-}{2}\le\frac{t}{[\zeta+\min(u,\zeta/3)]^2},\quad
 \E\binom{R_-}{3}\le\frac{t\sqrt{A^2+3t}}{(2\zeta-1/2)^3}.
\end{equation}
These are the crossing-moment bounds of Lemma~\ref{ms:crossing-curvature}
with $A_U=A$, $V_k=W_k$, $|u|\le1/2$ and the increment lower bound
$\zeta=1-d_0$. The lemma includes the strict outgoing crossing and the
fourth-moment estimate from independence and conditional centering.

\paragraph{The marked crossing.}
Let $D_i=U_{1,i}-U_{0,i}$, $\vartheta=\E(D_i)_+=\E(D_i)_-$,
$0<\eta\le1$, and $R_f=\E\{u+\eta-W_k\}$. The two deleted laws are
\[
 A_0\overset{d}=kU_1+(n-k-1)U_0,\qquad
 A_1\overset{d}=(k-1)U_1+(n-k)U_0.
\]
Then
\begin{align}
 \E R&\le u+\eta-R_f+J,\label{ms:fourier:cut49:shortloc:mark}\\
 J&=\frac{\vartheta}{\eta}\sum_{i\ge1}
 \{\Pp(A_0\le u+d_0-\zeta i)+\Pp(A_1>u+\eta+i-A)\}.\nonumber
\end{align}
Here are the pathwise details. Put $z=u-W_k$ and mark
\[
 Z_f=\sum_{i=1}^{n-k}(D_{k+i})_-\mathbf1_{\{W_k\le u-\zeta i\}},\quad
 Z_b=\sum_{i=1}^{k}(D_{k-i+1})_+\mathbf1_{\{W_k>u+\eta+i\}}.
\]
If $R>z+\eta$, the crossing activates the marks through $R$: for $R>0$
use $z\ge \zeta R$, while for $R<0$ use $z<R-\eta$. In either case
$R\le z+Z_f+Z_b$. Hence if $h=R-\lfloor z+\eta\rfloor>0$,
$Z_f+Z_b>h-1+\eta\ge\eta h$. Thus
$R\le\lfloor z+\eta\rfloor+(Z_f+Z_b)/\eta$ on every path.
On an active forward mark $U_0>U_1\ge-d_0$; deleting its pair implies
$A_0\le u+d_0-\zeta i$. On a backward mark delete $U_1\le A$ to obtain the
strict event for $A_1$. The remaining sum is independent of the entire
pair, so the expectations factor. Extending to all $i\ge1$ adds only
nonnegative bounds. Finally $|R|\le1+(|W_k|+|u|)/\zeta$ proves integrability
and permits expectations. No coupling has replaced either deleted marginal.

\paragraph{A finite bound for every mark.}
Fix $\lambda\in[L,U]$, $L>0$, and $t\le\omega_+\lambda$; put
$T_L=\omega_+L$, $T_U=\omega_+U$. On a phase subcell $u\in[v_l,v_h]$,
for $i=1,2$ let $b_1=\zeta i-d_0-v_h>0$ and $b_2=i-A+v_l>0$.
Lemma~\ref{ms:fourier:cut49:shortloc:tails} gives a bound $\omega_+\lambda K_i^{\rm ind}$,
where
\[
 K_i^{\rm ind}=\sum_{j=1}^2\min\left\{\frac1{b_j^2+T_L},
                          \frac{\overline P(A,T_U,b_j)}{T_L}\right\}.
\]
The denominator $T_L$ is a lower bound for the chosen upper envelope
$\omega_+\lambda$, not for any actual variance.
Alternatively couple $A_0=Z+U_1$, $A_1=Z+U_0$ with
$Z\overset d=(k-1)U_1+(n-k-1)U_0$. The two events project into
$Z\le-\alpha$ and $Z>\beta$, where
$\alpha=\zeta i-2d_0-v_h>0$, $\beta=i-A-d_0+v_l>0$.
For any $z\in(-\alpha,\beta)$ their sum is at most
$(\omega_+\lambda+z^2)/\min(z+\alpha,\beta-z)^2$.
Consequently a second kernel is
\[
 K_i^{\rm com}=\frac{1+z^2/T_L}{\min(z+\alpha,\beta-z)^2},\qquad
 z=\max\{-T_L/\beta,\min((\beta-\alpha)/2,T_L/\alpha)\}.
\]
This center is inside the interval; differentiating the two quadratic
ratios on either side of the midpoint proves its optimality, although
any interior rational center suffices. Take the smaller kernel on each
phase subcell and the maximum over their closed cover. Denote it $K_i$.
For $i\ge3$ use fourth moments and the convex midpoint integral for
$(cx-b)^{-4}$. On $u\in[u_l,u_h]$ this gives
\begin{equation}\label{ms:fourier:cut49:shortloc:K}
 K=K_1+K_2+\frac{A^2+3T_U}{3\zeta}
  \left\{\frac1{((5/2)\zeta-A-u_h)^3}
                 +\frac1{((5/2)\zeta-A+u_l)^3}\right\},\qquad
 J\le\frac{\vartheta}{\eta}\omega_+\lambda K.
\end{equation}
All denominators must be positive. Indeed, the tail fourth moment is at
most $A^2t+3t^2$, and integration from $5/2$ to infinity covers every
remaining integer, including the harmless added illegal counts.

If the same-law resource supplies $nd<G$ and
$G^2<c_vNL$, then
\begin{equation}\label{ms:fourier:cut49:shortloc:theta}
 \vartheta\le\frac{c_vG\lambda}{c_vn\lambda+G^2}.
\end{equation}
To prove this, for centered $X\ge-b$, $\E X^2\le V$ and $b^2\le V$,
the quadratic $b(X-V/b)^2/(b+V/b)^2$ majorizes $X_-$ on $[-b,\infty)$.
Its expectation is at most $bV/(V+b^2)$. Apply this to $D_i\ge-d$,
using $\operatorname{Var}(D_i)=s_0+s_1\le c_v\lambda/n$, and monotonicity
in $b\le\sqrt V$. The case $d=0$ is immediate. With $\eta=z\lambda$,
$z>0$, $zU\le1$, \eqref{ms:fourier:cut49:shortloc:K} becomes
\begin{equation}\label{ms:fourier:cut49:shortloc:Jfinal}
 \frac J\lambda\le\frac{\omega_+Kc_vG}{z(c_vN\lambda+G^2)}.
\end{equation}

\paragraph{The fractional-part gain.}
The following implication assumes the scalar majorization and its rational
constant enclosures stated below. Their finite arithmetic acceptance is a
separate condition.
If each centered summand has support span at most $0<b<1$ and the sum has
variance $t$, then, for any real $w$ and
$0<r_*\le2\sin^2(\pi b)/b^2$,
\begin{equation}\label{ms:fourier:cut49:shortloc:harmonic}
             \E\{w-W\}\ge\frac5{23}(1-e^{-r_*t}).
\end{equation}
Indeed $1-\cos(2\pi x)\le(23/5)x$ for $0\le x\le1$.
A finite proof of this sole scalar inequality is recorded below.
For an independent copy $Z'$ of one summand, concavity of sine on
$[0,\pi b]$ gives
\[
 |\E e^{2\pi iZ}|^2
 =1-2\E\sin^2\pi(Z-Z')
 \le e^{-4\sin^2(\pi b)\operatorname{Var}(Z)/b^2}.
\]
Multiply square roots over the independent summands and use periodicity.
This argument includes zero variances and integer collisions, and is
valid up to $b<1$ by Lemma~\ref{v9:harmonic:span}. The other count
denominator gates are unchanged.
Here the conditional spans are at most $A+d_0$, so this is the $b$ to use.

For a reproducible scalar proof, $x\ge10/23$ follows from $1-\cos\le2$.
On $[0,10/23]$ use the degree-ten alternating upper polynomial for
$1-\cos(2\pi x)$, with directed coefficients:
\[
 Q(x)=\sum_{j=1}^5\frac{(-1)^{j+1}(2\pi_jx)^{2j}}{(2j)!},\qquad
 \pi_j=\begin{cases}3.141592653589794,&j\text{ odd},\\
                       3.141592653589793,&j\text{ even}.\end{cases}
\]
Put
$P(x)=23/5-Q(x)/x$. On the eight intervals with left endpoint $a_j=5j/92$
and width $w=5/92$, its degree-nine Bernstein coefficients are
\[
 b_{jk}=\sum_{i=0}^k\frac{\binom ki}{\binom9i}
           \sum_{h=i}^9[x^h]P(x)\binom hi a_j^{h-i}w^i,
 \qquad 0\le j<8,\quad0\le k\le9.
\]
All 80 rational coefficients exceed $4492/100000$. This finite
comparison is reproduced by the independent rational check accompanying
the manuscript; together with the displayed formula it proves the scalar
majorization. Expansion in powers of the unit-interval coordinate
proves the formula; Bernstein weights are nonnegative and sum to one.
The alternating remainder and the stated rational bounds on $\pi$ follow
from the same elementary series above. This completely specifies the
scalar comparison. To obtain a rational rate throughout $0<b<1$, put
$y=\min(b,1-b)$ and $x=3.14159y$. Since $0<y\le1/2$,
$\sin(\pi b)=\sin(\pi y)\ge\sin x$, and the alternating bound
\[
 0<x-x^3/6+x^5/120-x^7/5040\le\sin x
\]
is valid. Positivity follows from $x^2<3$ for the first two terms
and $x^2<42$ for the last pair. Squaring this positive rational lower
bound therefore gives a valid rate below $2\sin^2(\pi b)/b^2$.
Folding at $1/2$ is essential: decreasing the sine argument would not be
monotone on the whole unfurled interval.

Set $r=r_*\omega_-$, require $rU<3$, and define
\begin{equation}\label{ms:fourier:cut49:shortloc:gain}
 g=\frac5{23}\left(\frac{r^3L^2}{6}-\frac{r^4L^3}{24}
                     +\frac{r^5L^4}{120}-\frac{r^6L^5}{720}\right).
\end{equation}
The sixth-order Taylor lower bound for $1-e^{-r\lambda}$ gives
$R_f/\lambda\ge5r/23-5r^2\lambda/46+g$. The discarded higher-order
part is increasing as a function of $\lambda$ for $r\lambda<3$:
its derivative has the sign of $48-18x+(24/5)x^2-x^3>0$ on $[0,3]$.
For instance this polynomial has positive degree-three Bernstein
coefficients $48,30,132/5,51/5$ there. Thus freezing it at $L$ is valid.

\paragraph{One final expression.}
On a signed phase cell find an interval $[x_l,x_h]$ containing $H/\ell$
by its four endpoint quotients, using positive bounds for $\ell$.
Let
\[
 D=\max(|x_l-u_h|,|x_h-u_l|),\quad
 D_+=\max(0,x_h-u_l),\quad a=\max(|u_l|,|u_h|),\quad b_\chi=p_+-x_l.
\]
The same formula allows coarser upper bounds $D,D_+,b_\chi$, as used by
the central certificate. Check $D\le\Delta$ and $b_\chi+u_h<2$.
With $\omega_\pm=1\pm c_vD/N$, the curvature coefficient is
\begin{equation}\label{ms:fourier:cut49:shortloc:V}
 V_0=\frac{\omega_+}{v_*}\left[
 \sup_{u_l\le u\le u_h}\left\{
 \frac{(b_\chi+u-1)_+}{(\zeta-u)^2}
 +\frac{(b_\chi+u)_+}{[\zeta+\min(u,\zeta/3)]^2}\right\}
 +\frac{\sqrt{A^2+3T_U}}{(2\zeta-1/2)^3}\right].
\end{equation}
On a rational phase subcell, both numerators are bounded at the right
endpoint; the first denominator is smallest there, and the second is
smallest at the left endpoint. This proves the supremum bound, including
both kinks. Equations \eqref{ms:fourier:cut49:shortloc:cutoff}--\eqref{ms:fourier:cut49:shortloc:crossingmoments},
$\sqrt n m_B\le1$, and \eqref{ms:fourier:cut49:shortloc:master} now give the whole-count
witness with
\begin{align}\label{ms:fourier:cut49:shortloc:BI}
 B(\lambda)&=z+\frac{\omega_+Kc_vG}{z(c_vN\lambda+G^2)}
                      -\frac{5r}{23}+\frac{5r^2\lambda}{46}-g,\\
 I(\lambda)&=V_0+aB_1+\frac{aP_1}{\lambda}+P_*D_+-m_cQ_*.
 \nonumber
\end{align}
We used $|u-f(u)|\le a(B_1\lambda+P_1)/n$ and $\sqrt n\pi_k\le1$.
Every constant is fixed before $\lambda$ varies. Both $B$ and
$mB+I/N$ are convex on $\lambda>0$; $I$ is decreasing. Thus their two
endpoint gates $B<0$, $I>0$, $mB+I/N<-0.001$ prove the entire closed noise
cell. Thus this argument supplies the analytic implication for each of the three
whole-count estimates, conditional on its own complete cover and endpoint gates.


\paragraph{The zero-noise endpoint and the remaining composition}
\label{ms:fourier:zero-noise}
If $\lambda=0$, then $p,q>0$ and conditional centering imply
$U_0=U_1=0$ almost surely. The actual law is its Bernoulli comparison law.
The zero-noise argument in Section~\ref{ms:counts} places an attained
positive closed-CDF maximum at a binomial atom, so $u=0$ and $C=C_B$.
The external all-probability binomial theorem of
\cite[Theorem~1]{Schulz2016}, when applied with its stated hypotheses,
then supplies the sharp Bernoulli upper bound. An equality-inclusive
exclusion requires a separately justified strict finite-index version. For positive noise, the local
conclusion requires every parameter cell to satisfy one of the displayed
strict tests. No selection of historical table rows, completeness of a
finite cover, or fresh acceptance of those tests is asserted here.

The mathematical composition therefore still requires: the actual contact
witness; already acquired common-law support, moments and signed resources;
any additional concentration hypotheses used in acquiring confinement; valid normalization and positivity
of every divisor; a complete cover of all signs, phases and noise values;
and acceptance of the stated finite gates on that cover. No step replaces
an $n-2$ convolution by an $n-1$ convolution or infers a uniform statement
from a fixed-law asymptotic limit.
